\chapter{Erectile Dysfunction}

\section*{Definition and Overview}
Erectile dysfunction is the persistent difficulty in achieving or maintaining an erection sufficient for satisfactory sexual activity. It is common, affecting a significant proportion of men, and becomes more frequent with age. Erectile dysfunction is often a sign of underlying health problems, particularly heart and blood vessel disease, and it can also be caused by diabetes, medicines, hormonal problems, and psychological factors. Effective treatments are available, and for many men erectile dysfunction is both treatable and a valuable early warning sign for other health conditions.

\section*{Epidemiology}
Erectile dysfunction affects a large proportion of Australian men, with prevalence increasing with age: it is estimated to affect about half of men over 50 and a higher proportion of older men, although it can occur at any age. It is under-reported because many men do not seek help, despite effective treatment. The condition is strongly linked to cardiovascular disease, diabetes, obesity, smoking, and mental health problems.

\section*{Aetiology and Pathophysiology}
An erection depends on healthy blood vessels, nerves, hormones, and psychological wellbeing.
\begin{itemize}
    \item \textbf{Blood vessel disease:} the most common cause; atherosclerosis narrows the arteries that fill the penis, and erectile dysfunction is often the first sign of cardiovascular disease
    \item \textbf{Diabetes:} damages the nerves and blood vessels
    \item \textbf{Neurological causes:} spinal cord injury, multiple sclerosis, and prostate surgery
    \item \textbf{Hormonal causes:} low testosterone and thyroid problems
    \item \textbf{Medicines:} some blood pressure medicines, antidepressants, and prostate medicines
    \item \textbf{Psychological factors:} stress, anxiety, depression, and relationship problems
    \item \textbf{Lifestyle:} smoking, excess alcohol, obesity, and inactivity
    \item \textbf{Mechanism:} an erection requires the arteries to dilate and the veins to close, all controlled by nerve signals and hormones; any interruption to this process causes erectile dysfunction
\end{itemize}

\section*{Clinical Features}
The presentation varies between men.
\begin{itemize}
    \item Difficulty achieving or maintaining an erection
    \item Reduced firmness of erections
    \item Reduced interest in sex, which may suggest a hormonal or psychological cause
    \item The problem may be constant or situational
    \item Morning erections may be present in psychological causes and absent in physical causes
    \item Associated symptoms of the underlying condition, such as fatigue, reduced libido, or urinary symptoms
\end{itemize}

\section*{Differential Diagnosis}
Other conditions affect sexual function.
\begin{itemize}
    \item Premature ejaculation and other ejaculatory disorders
    \item Low libido from depression, hormonal problems, or relationship issues
    \item Peyronie's disease, with curvature of the penis
    \item Delayed ejaculation
    \item The distinction matters because the treatments differ
\end{itemize}

\section*{When to Seek Medical Care}
All men with erectile dysfunction should discuss it with a GP, both because it is treatable and because it may indicate cardiovascular risk. Men with sudden erectile dysfunction after injury or surgery, or with associated symptoms, need medical assessment.

\textbf{Call 000 (or local emergency services) for:}
\begin{itemize}
    \item Sudden loss of erections with back pain and leg weakness
    \item Erection lasting more than four hours after taking treatment (priapism)
    \item Chest pain or other signs of a heart attack
    \item Sudden severe pelvic or testicular pain
\end{itemize}

\section*{Diagnosis and Investigations}
Assessment aims to identify the cause and any underlying health risks.
\begin{itemize}
    \item \textbf{History:} the pattern of the problem, medicines, lifestyle, and sexual history
    \item \textbf{Examination:} blood pressure, pulse, and examination of the genital area
    \item \textbf{Blood tests:} glucose, cholesterol, testosterone, and thyroid function
    \item \textbf{Cardiovascular risk assessment:} because erectile dysfunction predicts heart disease
    \item \textbf{Specialist referral:} to a urologist for complex or persistent cases
\end{itemize}

\section*{Management}
Treatment is effective and tailored to the cause.
\begin{itemize}
    \item \textbf{Lifestyle changes:} exercise, weight loss, stopping smoking, and reducing alcohol
    \item \textbf{Treating the cause:} improving diabetes control, blood pressure, and cholesterol; changing medicines that cause the problem
    \item \textbf{Phosphodiesterase inhibitors:} medicines such as sildenafil and tadalafil, taken before sex, which are highly effective
    \item \textbf{Other treatments:} vacuum devices, injections, and penile implants for men who do not respond to tablets
    \item \textbf{Hormone therapy:} testosterone for men with deficiency
    \item \textbf{Psychological support:} counselling and sex therapy for psychological causes
    \item \textbf{Involving the partner:} relationship support improves outcomes
\end{itemize}

\section*{Complications}
\begin{itemize}
    \item The psychological effects: anxiety, low self-esteem, and depression
    \item Relationship strain
    \item Undiagnosed cardiovascular disease, of which erectile dysfunction is a warning sign
    \item Side effects of treatments
    \item Rarely, priapism, a prolonged painful erection requiring urgent care
\end{itemize}

\section*{Prognosis and Outcomes}
Erectile dysfunction is highly treatable, and most men respond to oral medicines or other treatments. Addressing the underlying risk factors improves both sexual function and general health. For men whose erectile dysfunction is caused by vascular disease, treatment is an opportunity to reduce the risk of heart attack and stroke.

\section*{Prevention and Self-Care}
\begin{itemize}
    \item Exercise regularly and maintain a healthy weight
    \item Do not smoke, and limit alcohol
    \item Manage diabetes, blood pressure, and cholesterol
    \item Discuss erectile dysfunction openly with a doctor
    \item Manage stress and address relationship concerns
    \item Take treatments as directed and report side effects
    \item Seek urgent care for a prolonged erection
\end{itemize}

\section*{Sources and Further Reading}
The following reputable sources provide further information for healthcare students and patients seeking reliable, current advice about erectile dysfunction.
\begin{itemize}
    \item Healthy Male Australia. \textit{Erectile dysfunction information.} https://www.healthymale.org.au/
    \item Healthdirect Australia. \textit{Erectile dysfunction.} https://www.healthdirect.gov.au/erectile-dysfunction
    \item Andrology Australia. \textit{Male sexual health resources.} https://www.andrologyaustralia.org/
    \item NHS. \textit{Erectile dysfunction (impotence).} https://www.nhs.uk/conditions/erectile-dysfunction/
    \item Mayo Clinic. \textit{Erectile dysfunction: symptoms and treatment.} https://www.mayoclinic.org/diseases-conditions/erectile-dysfunction/
    \item Urology Care Foundation. \textit{Erectile dysfunction patient information.} https://www.urologyhealth.org/
\end{itemize}
